\documentclass[10pt,letterpaper]{article}

\usepackage[letterpaper,margin=0.72in,headheight=20pt,footskip=28pt]{geometry}
\usepackage[T1]{fontenc}
\usepackage[utf8]{inputenc}
\usepackage{lmodern}
\usepackage{microtype}
\usepackage[dvipsnames,table]{xcolor}
\usepackage{booktabs}
\usepackage{tabularx}
\usepackage{longtable}
\usepackage{array}
\usepackage{enumitem}
\usepackage{fancyhdr}
\usepackage{titlesec}
\usepackage[most]{tcolorbox}
\usepackage{listings}
\usepackage{hyperref}
\usepackage{lastpage}
\usepackage{ragged2e}

% ---------- Palette ----------
\definecolor{Navy}{HTML}{17324D}
\definecolor{Blue}{HTML}{246B91}
\definecolor{Teal}{HTML}{1C7C78}
\definecolor{Gold}{HTML}{C7902F}
\definecolor{Slate}{HTML}{52606D}
\definecolor{Ink}{HTML}{17212B}
\definecolor{Mist}{HTML}{EEF4F7}
\definecolor{PaleBlue}{HTML}{F4F8FB}
\definecolor{PaleGold}{HTML}{FFF8E8}
\definecolor{Rule}{HTML}{D7E0E7}

\color{Ink}
\hypersetup{
  colorlinks=true,
  linkcolor=Navy,
  urlcolor=Blue,
  citecolor=Teal,
  pdftitle={Self-Learning Curriculum: Documenting Software Architectures},
  pdfsubject={A 10-week Views and Beyond architecture documentation curriculum},
  pdfauthor={Prepared as a professional self-study handout}
}
\urlstyle{same}

% ---------- Typography ----------
\titleformat{\section}
  {\Large\bfseries\color{Navy}}
  {}{0pt}{}
  [\vspace{-0.35em}\color{Rule}\titlerule]
\titleformat{\subsection}
  {\large\bfseries\color{Blue}}
  {}{0pt}{}
\titleformat{\subsubsection}
  {\normalsize\bfseries\color{Navy}}
  {}{0pt}{}
\titlespacing*{\section}{0pt}{1.35em}{0.65em}
\titlespacing*{\subsection}{0pt}{1.0em}{0.35em}
\titlespacing*{\subsubsection}{0pt}{0.7em}{0.2em}

\setlength{\parindent}{0pt}
\setlength{\parskip}{0.42em}
\setlist[itemize]{leftmargin=1.25em,itemsep=0.16em,topsep=0.25em}
\setlist[enumerate]{leftmargin=1.35em,itemsep=0.16em,topsep=0.25em}
\renewcommand{\arraystretch}{1.18}

% ---------- Header / footer ----------
\pagestyle{fancy}
\fancyhf{}
\fancyhead[L]{\small\color{Slate}\textbf{SOFTWARE ARCHITECTURE} \textbar{} SELF-STUDY CURRICULUM}
\fancyhead[R]{\small\color{Slate}Views and Beyond}
\fancyfoot[L]{\small\color{Slate}Professional self-learning handout}
\fancyfoot[R]{\small\color{Slate}Page \thepage\ of \pageref*{LastPage}}
\renewcommand{\headrulewidth}{0.4pt}
\renewcommand{\footrulewidth}{0pt}
\renewcommand{\headrule}{\hbox to\headwidth{\color{Rule}\leaders\hrule height \headrulewidth\hfill}}

% ---------- Table helpers ----------
\newcolumntype{Y}{>{\RaggedRight\arraybackslash}X}
\newcolumntype{P}[1]{>{\RaggedRight\arraybackslash}p{#1}}

% ---------- Boxes ----------
\newtcolorbox{keybox}[1][]{
  enhanced,
  colback=Mist,
  colframe=Teal,
  boxrule=0.7pt,
  arc=2pt,
  left=8pt,right=8pt,top=6pt,bottom=6pt,
  #1
}
\newtcolorbox{deliverablebox}{
  enhanced,
  colback=PaleBlue,
  colframe=Blue,
  boxrule=0.6pt,
  arc=2pt,
  left=8pt,right=8pt,top=5pt,bottom=5pt,
  title=\textbf{Weekly deliverable},
  coltitle=Navy,
  fonttitle=\bfseries
}
\newtcolorbox{masterybox}{
  enhanced,
  colback=PaleGold,
  colframe=Gold,
  boxrule=0.6pt,
  arc=2pt,
  left=8pt,right=8pt,top=5pt,bottom=5pt,
  title=\textbf{Mastery check},
  coltitle=Navy,
  fonttitle=\bfseries
}

\lstset{
  basicstyle=\ttfamily\small,
  backgroundcolor=\color{PaleBlue},
  frame=single,
  rulecolor=\color{Rule},
  framesep=5pt,
  xleftmargin=3pt,
  xrightmargin=3pt,
  breaklines=true,
  showstringspaces=false,
  columns=fullflexible
}

\begin{document}

% ============================================================================
% TITLE
% ============================================================================
\thispagestyle{empty}
\begin{tcolorbox}[
  enhanced,
  colback=Navy,
  colframe=Navy,
  boxrule=0pt,
  arc=3pt,
  left=18pt,right=18pt,top=18pt,bottom=18pt]
{\color{white}\small\bfseries CARNEGIE MELLON SEI-ALIGNED SELF-STUDY}\par
\vspace{0.35em}
{\color{white}\fontsize{25}{29}\selectfont\bfseries Self-Learning Curriculum:\par
Documenting Software Architectures}\par
\vspace{0.55em}
{\color{white!82}\large A 10-week practical program based on the SEI Views and Beyond approach}
\end{tcolorbox}

\vspace{1.1em}
\begin{keybox}
\textbf{Purpose.} Build the ability to produce a comprehensive, stakeholder-oriented software architecture description rather than a collection of disconnected diagrams. The curriculum follows the objectives and topic areas of Carnegie Mellon University's Software Engineering Institute (SEI) course \emph{Documenting Software Architectures} and uses \emph{Documenting Software Architectures: Views and Beyond, 2nd Edition} as its principal paid reference.
\end{keybox}

\vspace{1.0em}
\begin{tabularx}{\textwidth}{@{}P{0.23\textwidth}Y@{}}
\toprule
\rowcolor{Mist}\textbf{Program element} & \textbf{Recommendation} \\
\midrule
Duration & 10 weeks \\
Weekly workload & 4--6 hours \\
Total effort & Approximately 50 hours \\
Mastery threshold & 80\% on weekly self-assessments \\
Primary method & SEI Views and Beyond (V\&B) \\
Primary case study & AI-powered SAST scanner / security analysis platform \\
Capstone & Complete, reviewable software architecture documentation package \\
Standards alignment & ISO/IEC/IEEE 42010:2022 concepts, plus UML 2.5.1 notation where useful \\
\bottomrule
\end{tabularx}

\vfill
\begin{tcolorbox}[
  enhanced,
  colback=PaleGold,
  colframe=Gold,
  boxrule=0.7pt,
  arc=2pt,
  left=10pt,right=10pt,top=8pt,bottom=8pt]
\textbf{Recommended learning loop}\par
\centering\large
\textbf{Read} $\rightarrow$ \textbf{Model a real system} $\rightarrow$ \textbf{Document} $\rightarrow$ \textbf{Review} $\rightarrow$ \textbf{Revise}
\end{tcolorbox}

\vspace{0.9em}
{\small\color{Slate}
Prepared August 2026. Resource availability and listed prices are snapshots and may change. Official source links are provided in the resource guide.}
\newpage

% ============================================================================
\section{Program Outcomes}

By the end of the program, the learner should be able to:

\begin{enumerate}
  \item Apply principles of sound technical documentation from the reader's point of view.
  \item Identify architecture stakeholders, their concerns, and the intended uses of documentation.
  \item Distinguish \textbf{viewtypes}, \textbf{styles}, \textbf{structures}, and concrete \textbf{views}.
  \item Create useful module, component-and-connector, and allocation views.
  \item Document context, refinement, variability, interfaces, behavior, and architectural rationale.
  \item Choose views systematically based on stakeholder needs instead of producing diagrams by habit.
  \item Relate Views and Beyond documentation to current ISO/IEC/IEEE 42010:2022 concepts.
  \item Assemble, review, and maintain a coherent architecture documentation package.
\end{enumerate}

\begin{keybox}
\textbf{Capstone philosophy.} The measure of success is not finishing the readings. It is being able to take a nontrivial software system---ideally a real repository---and produce a stakeholder-driven, multi-view, internally consistent architecture description that another engineer can use.
\end{keybox}

\subsection{10-Week Roadmap}

\begin{longtable}{@{}P{0.08\textwidth}P{0.35\textwidth}P{0.49\textwidth}@{}}
\toprule
\rowcolor{Navy}\color{white}\textbf{Week} & \color{white}\textbf{Theme} & \color{white}\textbf{Primary output} \\
\midrule
\endhead
1 & Architecture documentation fundamentals & Stakeholder and concern inventory \\
2 & Views, viewtypes, and styles & Architecture view catalog \\
3 & Module views & Static / code-oriented architecture view \\
4 & Component-and-connector views & Runtime architecture view \\
5 & Allocation views & Deployment, implementation, and work-assignment views \\
6 & Advanced documentation concepts & System context and architectural rationale \\
7 & Software interfaces & Interface specifications \\
8 & Architectural behavior & Behavioral models and finding lifecycle \\
9 & Selecting views + ISO 42010 & Stakeholder--concern--viewpoint--view matrix \\
10 & Package construction and review & Complete capstone architecture description \\
\bottomrule
\end{longtable}

% ============================================================================
\section{Week 1 --- Principles of Sound Architecture Documentation}

\textbf{Goal:} Understand architecture documentation as a communication and engineering artifact, not merely a set of pictures.

\subsection{Learn}
\begin{itemize}
  \item Software architecture versus architecture documentation.
  \item Stakeholders, concerns, intended uses, structures, views, and styles.
  \item Reader-centered technical documentation.
  \item SEI's rules for sound documentation, including standard organization and precise notation.
\end{itemize}

\subsection{Read / Watch}
\begin{itemize}
  \item \href{https://www.informit.com/articles/article.aspx?p=1641654}{\emph{Prologue: Software Architectures and Documentation}} --- free publisher excerpt.
  \item \href{https://www.informit.com/articles/article.aspx?p=1641654&seqNum=5}{\emph{Seven Rules for Sound Documentation}} --- free publisher excerpt.
  \item Optional: V\&B 2nd ed., Prologue.
\end{itemize}

\subsection{Practice}
Identify representative stakeholders for the SAST scanner, such as application security engineers, developers, security architects, DevOps/platform engineers, repository administrators, security managers, and auditors. For each, record the role, concerns, questions that the documentation must answer, and candidate views.

\begin{deliverablebox}
\texttt{01-stakeholders-and-concerns.md} --- stakeholder inventory, concerns, documentation questions, and initial view candidates.
\end{deliverablebox}

\begin{masterybox}
Explain without notes: \emph{Why is a single architecture diagram insufficient as architecture documentation?} Score yourself on correctness, completeness, clarity, and use of stakeholder concerns. Target: \textbf{80\% or better}.
\end{masterybox}

% ============================================================================
\section{Week 2 --- Viewtypes, Styles, and Views}

\textbf{Goal:} Master the conceptual vocabulary that underpins Views and Beyond.

\subsection{Learn the Three Fundamental V\&B Viewtypes}
\begin{tabularx}{\textwidth}{@{}P{0.26\textwidth}Y@{}}
\toprule
\rowcolor{Mist}\textbf{Viewtype} & \textbf{Question it answers} \\
\midrule
Module & How is the software organized as implementation units? \\
Component-and-Connector (C\&C) & How does the system operate and interact at runtime? \\
Allocation & How is software mapped to nonsoftware structures in its environment? \\
\bottomrule
\end{tabularx}

\subsection{Study}
Distinguish structure, viewtype, style, and view. Treat diagrams as representations of architectural information; the architecture is not the diagram itself. Review the \href{https://www.sei.cmu.edu/library/views-and-beyond-collection/}{SEI Views and Beyond Collection} and its explanation of views as the fundamental organizing principle.

\subsection{Practice}
Classify 15 architecture diagrams from documentation or real systems as module, C\&C, allocation, hybrid, or non-architectural. Then classify the existing diagrams for the capstone system and note ambiguities.

\begin{deliverablebox}
\texttt{02-view-catalog.md} --- view inventory, classifications, intended stakeholders, and purpose of each view.
\end{deliverablebox}

\begin{masterybox}
Given an unfamiliar diagram, identify the kinds of elements and relations being represented and justify its viewtype. Target: \textbf{80\% correct across 10 examples}.
\end{masterybox}

% ============================================================================
\section{Week 3 --- Module Views}

\textbf{Goal:} Describe the static implementation organization of a system.

\subsection{Read}
\begin{itemize}
  \item V\&B 2nd ed., Chapter 1: \emph{Module Views}.
  \item V\&B 2nd ed., Chapter 2: \emph{A Tour of Some Module Styles}.
\end{itemize}

\subsection{Study}
\begin{itemize}
  \item Decomposition, uses, layered, and generalization styles.
  \item Module responsibilities and properties.
  \item Relationships and dependencies.
  \item Interfaces and constraints.
\end{itemize}

\subsection{Practice}
Create a top-level module decomposition for the SAST scanner. Candidate responsibilities include CLI, repository analysis, language analyzers, taint analysis, exploitability classification, finding management, SARIF generation, reporting, GHAS integration, and remediation tracking. Then document actual dependency relationships; do not assume the hierarchy alone captures them.

\begin{deliverablebox}
\texttt{03-module-view.md} --- primary presentation, element descriptions, responsibilities, relationships, constraints, and rationale.
\end{deliverablebox}

\begin{masterybox}
For every module shown, answer: \emph{What is it responsible for? What does it depend on? What depends on it? What architectural rule constrains it?}
\end{masterybox}

% ============================================================================
\section{Week 4 --- Component-and-Connector Views}

\textbf{Goal:} Describe runtime elements, connectors, communication paths, and interaction semantics.

\subsection{Read}
\begin{itemize}
  \item V\&B 2nd ed., Chapter 3: \emph{Component-and-Connector Views}.
  \item V\&B 2nd ed., Chapter 4: \emph{A Tour of Some Component-and-Connector Styles}.
\end{itemize}

\subsection{Study}
Client-server, service-oriented, publish-subscribe, shared-data, pipe-and-filter, peer-to-peer, and communicating-process styles. Keep the distinction explicit: \textbf{module views describe implementation structure; C\&C views describe runtime structure}.

\subsection{Practice}
Model a complete scan execution from invocation through analysis, finding normalization, SARIF/report generation, and GitHub/GHAS integration. For each connector document direction, information transferred, protocol/API, synchronization, failure behavior, trust boundary, and data classification where relevant.

\begin{deliverablebox}
\texttt{04-runtime-view.md} --- runtime components, connector semantics, runtime responsibilities, and relevant quality/security notes.
\end{deliverablebox}

\begin{masterybox}
Demonstrate that every line or arrow in the primary presentation has explicit semantics. A reader should not need tribal knowledge to determine what a connector means.
\end{masterybox}

% ============================================================================
\section{Week 5 --- Allocation Views}

\textbf{Goal:} Connect the software architecture to infrastructure, source/repository organization, and development responsibility.

\subsection{Read}
\begin{itemize}
  \item V\&B 2nd ed., Chapter 5: \emph{Allocation Views and a Tour of Some Allocation Styles}.
\end{itemize}

\subsection{Study}
\begin{tabularx}{\textwidth}{@{}P{0.28\textwidth}Y@{}}
\toprule
\rowcolor{Mist}\textbf{Allocation style} & \textbf{Mapping} \\
\midrule
Deployment & Software $\rightarrow$ execution infrastructure \\
Implementation & Software $\rightarrow$ repositories, directories, packages, build/configuration structures \\
Work assignment & Software $\rightarrow$ teams or engineering ownership \\
\bottomrule
\end{tabularx}

\subsection{Practice}
Produce all three allocations for the capstone. Give extra attention to the implementation view because it creates a bridge between conceptual architecture and the physical repository structure---a prerequisite for repository-aware architecture analysis.

\begin{deliverablebox}
\texttt{05-allocation-views.md} --- deployment, implementation, and work-assignment views with allocation rules and rationale.
\end{deliverablebox}

\begin{masterybox}
Trace one important capability from source location to runtime deployment to engineering ownership. Identify any gaps or ambiguous ownership.
\end{masterybox}

% ============================================================================
\section{Week 6 --- Context, Refinement, Variability, and Rationale}

\textbf{Goal:} Capture the information that makes individual views useful and makes the architecture understandable as a whole.

\subsection{Read}
\begin{itemize}
  \item V\&B 2nd ed., Chapter 6: \emph{Beyond the Basics}.
  \item Revisit the \href{https://www.sei.cmu.edu/training/documenting-software-architectures/}{SEI course objectives and topics} for refinement, context diagrams, variability, interfaces, and behavior.
\end{itemize}

\subsection{Study}
\begin{itemize}
  \item System boundary and external context.
  \item Information chunking and refinement.
  \item View packets and relationships between views.
  \item Variability and dynamic architectures.
  \item Architectural background, constraints, and rationale.
\end{itemize}

\subsection{Practice}
Create a SAST system context including GitHub, repositories, GHAS/code scanning, AI/model services, CI/CD, developers, AppSec analysts, and report consumers as applicable. Start an Architecture Decision Record (ADR) collection for significant decisions.

\begin{lstlisting}
Decision
Context
Alternatives considered
Selected approach
Rationale
Consequences
Status
\end{lstlisting}

\begin{deliverablebox}
\texttt{06-system-context.md} plus \texttt{decisions/} containing initial ADRs.
\end{deliverablebox}

\begin{masterybox}
Verify that the system boundary, every external dependency, and the rationale for at least three significant architecture decisions can be understood without interviewing the author.
\end{masterybox}

% ============================================================================
\section{Week 7 --- Documenting Software Interfaces}

\textbf{Goal:} Specify architectural interfaces with enough semantics to support implementation, integration, analysis, and maintenance.

\subsection{Read}
\begin{itemize}
  \item V\&B 2nd ed., Chapter 7: \emph{Documenting Software Interfaces}.
\end{itemize}

\subsection{Document More Than an Endpoint}
For important interfaces, consider identity, provided and required operations, inputs, outputs, preconditions, postconditions, data types, errors, protocols, quality requirements, versioning, security requirements, and constraints.

\subsection{Practice}
Fully specify the scanner's integration with GitHub code scanning/SARIF, then specify one significant internal interface. Include trust and authentication properties where the interface crosses a security boundary.

\begin{deliverablebox}
\texttt{07-interface-specifications.md} --- at least two complete interface specifications.
\end{deliverablebox}

\begin{masterybox}
Ask whether a separate engineering team could implement a compatible peer using only the interface documentation. Record every unanswered question as a defect to correct.
\end{masterybox}

% ============================================================================
\section{Week 8 --- Documenting Architectural Behavior}

\textbf{Goal:} Represent how architectural elements and the system behave over time.

\subsection{Read}
\begin{itemize}
  \item V\&B 2nd ed., Chapter 8: \emph{Documenting Behavior}.
  \item \href{https://www.omg.org/spec/UML/}{OMG UML 2.5.1 specification} as a free notation reference where UML is appropriate.
\end{itemize}

\subsection{Study}
Select notation based on the question being answered. Practice sequence diagrams, activity diagrams, state machines, and precise narrative behavior specifications. Avoid diagrams whose semantics are undefined.

\subsection{Practice}
Model at least three scenarios: successful scan, failed scan, and finding detected $\rightarrow$ SARIF produced $\rightarrow$ GHAS uploaded. Also model the lifecycle of a finding from detection through triage, remediation, verification, and closure.

\begin{deliverablebox}
\texttt{08-behavior.md} --- behavioral models, legends/keys, assumptions, and narrative explanations.
\end{deliverablebox}

\begin{masterybox}
For each behavior model, identify its initiating stimulus, participants, important state changes, failure paths, and termination condition.
\end{masterybox}

% ============================================================================
\section{Week 9 --- Choosing Views and Applying ISO/IEC/IEEE 42010}

\textbf{Goal:} Select architecture views because stakeholders need them, not because a framework provides a checklist.

\subsection{Read}
\begin{itemize}
  \item V\&B 2nd ed., Chapter 9: \emph{Choosing the Views}.
  \item \href{https://www.sei.cmu.edu/library/comparing-the-seis-views-and-beyond-approach-for-documenting-software-architectures-with-ansi-ieee-1471-2000/}{SEI: Comparing Views and Beyond with ANSI/IEEE 1471-2000} --- free historical crosswalk.
  \item \href{https://www.iso.org/standard/74393.html}{ISO/IEC/IEEE 42010:2022} --- official current standard page and free preview.
\end{itemize}

\begin{keybox}
\textbf{Standards note.} The SEI comparison report discusses IEEE 1471-2000, which is historically important to the multi-view architecture-description model. For present-day standards work, use the current \textbf{ISO/IEC/IEEE 42010:2022, Edition 2}. Do not treat the 2005 comparison report as a substitute for the current standard.
\end{keybox}

\subsection{Practice}
Construct a stakeholder-to-view selection matrix. A starter example follows.

\begin{tabularx}{\textwidth}{@{}P{0.18\textwidth}P{0.22\textwidth}P{0.25\textwidth}Y@{}}
\toprule
\rowcolor{Mist}\textbf{Stakeholder} & \textbf{Concern} & \textbf{Viewpoint} & \textbf{View} \\
\midrule
Developer & Code organization & Module & Decomposition \\
AppSec & Trust boundaries & Runtime/security & C\&C \\
DevOps & Execution environment & Allocation & Deployment \\
Architect & System interaction & C\&C & Runtime \\
Auditor & Security decisions & Cross-view & Rationale / ADRs \\
\bottomrule
\end{tabularx}

\begin{deliverablebox}
\texttt{09-viewpoint-matrix.md} --- stakeholders, concerns, viewpoints/model kinds, selected views, and justification for inclusion.
\end{deliverablebox}

\begin{masterybox}
Defend every included view: identify who uses it, the concern it addresses, and what decision or task it enables. Remove views that have no clear consumer or purpose.
\end{masterybox}

% ============================================================================
\section{Week 10 --- Build and Review the Architecture Package}

\textbf{Goal:} Integrate the prior nine weeks into a coherent, usable, reviewable architecture description.

\subsection{Read}
\begin{itemize}
  \item V\&B 2nd ed., Chapter 10: \emph{Building the Documentation Package}.
  \item V\&B 2nd ed., Chapter 11: \emph{Reviewing an Architecture Document}.
  \item V\&B 2nd ed., Epilogue: using V\&B with other approaches.
  \item \href{https://www.sei.cmu.edu/library/views-and-beyond-documentation-template/}{SEI Views and Beyond Documentation Template} --- free.
  \item \href{https://www.sei.cmu.edu/library/a-structured-approach-for-reviewing-architecture-documentation/}{SEI: A Structured Approach for Reviewing Architecture Documentation} --- free.
\end{itemize}

\subsection{Recommended Capstone Repository Structure}
\begin{lstlisting}
architecture/
|-- README.md
|-- stakeholders/
|-- context/
|-- requirements/
|-- views/
|   |-- module/
|   |-- component-connector/
|   `-- allocation/
|-- interfaces/
|-- behavior/
|-- decisions/
|-- mappings/
|-- rationale/
|-- glossary/
`-- review/
\end{lstlisting}

\subsection{Stakeholder Review Questions}
\begin{itemize}
  \item Can developers implement or safely modify the system using the documentation?
  \item Can AppSec identify trust boundaries, sensitive flows, and security-relevant interfaces?
  \item Can operations determine where and how the system executes?
  \item Can testers identify architecturally significant behaviors and quality concerns?
  \item Can a new architect understand the architecture without substantial oral explanation?
  \item Can reviewers determine \emph{why} significant decisions were made?
  \item Can stakeholders find what they need without reading the entire package?
  \item Are diagram semantics explicit and are views mutually consistent?
\end{itemize}

\begin{deliverablebox}
\textbf{Capstone: Software Architecture Description --- AI-Powered SAST Scanner.} Include the complete package plus documented review findings and corrections.
\end{deliverablebox}

% ============================================================================
\section{Capstone Assessment Rubric}

Score each criterion from 0--4. A score of 3 means professionally usable; 4 means exceptionally clear, complete, and traceable. Target at least \textbf{80\% overall} and no zero in any category.

\begin{longtable}{@{}P{0.30\textwidth}P{0.58\textwidth}P{0.06\textwidth}@{}}
\toprule
\rowcolor{Navy}\color{white}\textbf{Criterion} & \color{white}\textbf{Evidence expected} & \color{white}\textbf{/4} \\
\midrule
\endhead
Stakeholder orientation & Stakeholders, concerns, uses, and navigation are explicit. & \\
View selection & Every view has a justified stakeholder purpose. & \\
Module documentation & Static structure, relations, properties, and responsibilities are clear. & \\
Runtime documentation & Runtime components and connector semantics are precise. & \\
Allocation documentation & Deployment/implementation mappings are usable and current. & \\
Interfaces & Architecturally significant interfaces include meaningful semantics and constraints. & \\
Behavior & Key success, failure, and state-transition scenarios are documented. & \\
Context and boundaries & System boundary, external dependencies, and relevant trust boundaries are clear. & \\
Rationale & Significant decisions, alternatives, consequences, and constraints are traceable. & \\
Cross-view consistency & Names, responsibilities, mappings, and relationships agree across views. & \\
Notation quality & Every diagram has an understandable legend/key and avoids ambiguous semantics. & \\
Maintainability & Organization, ownership, review status, and update paths are evident. & \\
\bottomrule
\end{longtable}

\textbf{Scoring:} 12 criteria $\times$ 4 points = 48 possible. Mastery target: \textbf{39/48 or higher} (approximately 81\%).

% ============================================================================
\section{Resource Guide --- Free}

The no-cost path is strong enough to complete the curriculum if purchasing the book is not currently desirable.

\begin{longtable}{@{}P{0.37\textwidth}P{0.48\textwidth}P{0.09\textwidth}@{}}
\toprule
\rowcolor{Navy}\color{white}\textbf{Resource} & \color{white}\textbf{Use} & \color{white}\textbf{Cost} \\
\midrule
\endhead
\href{https://www.sei.cmu.edu/library/views-and-beyond-collection/}{SEI Views and Beyond Collection} & Central collection of SEI V\&B publications and supporting material. & Free \\
\href{https://www.sei.cmu.edu/library/views-and-beyond-documentation-template/}{SEI V\&B Documentation Template} & Official starting template for a software architecture document. & Free \\
\href{https://www.informit.com/articles/article.aspx?p=1641654}{V\&B 2nd ed. Prologue} & Architecture/documentation fundamentals, views, and styles. & Free \\
\href{https://www.informit.com/articles/article.aspx?p=1641654&seqNum=5}{Seven Rules for Sound Documentation} & Objective checklist for documentation quality. & Free \\
\href{https://www.sei.cmu.edu/library/documenting-software-architectures-podcast/}{SEI Documenting Software Architectures Podcast} & Author-level overview of communicating architectures. & Free \\
\href{https://www.sei.cmu.edu/library/documenting-software-architectures-in-an-agile-world/}{Documenting Architectures in an Agile World} & Integrates architecture documentation with iterative/agile development. & Free \\
\href{https://www.sei.cmu.edu/library/comparing-the-seis-views-and-beyond-approach-for-documenting-software-architectures-with-ansi-ieee-1471-2000/}{V\&B vs. ANSI/IEEE 1471-2000} & Historical standards crosswalk and multi-view architecture-description concepts. & Free \\
\href{https://www.sei.cmu.edu/library/a-structured-approach-for-reviewing-architecture-documentation/}{Structured Architecture Documentation Review} & Stakeholder-centered method for capstone review. & Free \\
\href{https://www.omg.org/spec/UML/}{OMG UML 2.5.1} & Normative UML specification and notation reference. & Free \\
\href{https://www.iso.org/standard/74393.html}{ISO/IEC/IEEE 42010:2022 official page/preview} & Current architecture-description standard overview and sample. & Free preview \\
\bottomrule
\end{longtable}

% ============================================================================
\section{Resource Guide --- Paid}

\subsection{1. Core Book --- Recommended}
\textbf{Paul Clements et al., \emph{Documenting Software Architectures: Views and Beyond, 2nd Edition}.} Addison-Wesley Professional, ISBN 978-0-321-55268-6. The SEI course is based on this text.

\begin{tabularx}{\textwidth}{@{}P{0.27\textwidth}Y@{}}
\toprule
\rowcolor{Mist}\textbf{Option} & \textbf{Snapshot price / note} \\
\midrule
InformIT eBook & Approximately US\$61.59 \\
InformIT print book & Approximately US\$63.99 \\
Official page & \href{https://www.informit.com/store/documenting-software-architectures-views-and-beyond-9780321552686}{InformIT --- Views and Beyond, 2nd Edition} \\
\bottomrule
\end{tabularx}

\subsection{2. Official SEI eLearning --- Optional}
\href{https://www.sei.cmu.edu/training/documenting-software-architectures-elearning/}{SEI Documenting Software Architectures --- eLearning}. Listed at \textbf{US\$500}. SEI describes 13 instructional modules, quizzes, optional exercises with solution guidance, downloadable slides, a concluding assessment, six months of access, and 1.5 CEUs. The textbook is acquired separately.

\subsection{3. ISO/IEC/IEEE 42010:2022 --- Advanced / Professional}
\href{https://www.iso.org/standard/74393.html}{ISO/IEC/IEEE 42010:2022, Edition 2}. The ISO store listed the full standard at \textbf{CHF 204} at the time this handout was prepared. Begin with the free preview; purchase the full standard when formal architecture-description conformance is a project requirement.

\subsection{4. Live SEI Course --- Reference Point}
The \href{https://www.sei.cmu.edu/training/documenting-software-architectures/}{live SEI course} listed the following fees when this handout was prepared:

\begin{tabularx}{\textwidth}{@{}P{0.42\textwidth}Y@{}}
\toprule
\rowcolor{Mist}\textbf{Category} & \textbf{Listed fee} \\
\midrule
Government / Academic & US\$1,200 \\
Industry & US\$1,500 \\
International & US\$2,250 \\
\bottomrule
\end{tabularx}

\begin{keybox}
\textbf{Best-value recommendation.} For an experienced software/security engineer, start with the \textbf{V\&B book + free SEI collection + this 10-week capstone}. Add the official eLearning course if instructor-produced modules, exercises, CEUs, or structured assessment justify the additional cost. Purchase ISO 42010 when formal conformance becomes necessary.
\end{keybox}

% ============================================================================
\section{Suggested Weekly Routine}

\begin{tabularx}{\textwidth}{@{}P{0.22\textwidth}P{0.18\textwidth}Y@{}}
\toprule
\rowcolor{Navy}\color{white}\textbf{Activity} & \color{white}\textbf{Time} & \color{white}\textbf{Purpose} \\
\midrule
Read / watch & 1.5--2.0 h & Acquire concepts and vocabulary. \\
Model & 1.0--1.5 h & Apply the concepts to a real system. \\
Document & 1.0--1.5 h & Convert understanding into a durable architecture artifact. \\
Review / test & 0.5--1.0 h & Find ambiguity, inconsistency, missing rationale, and stakeholder gaps. \\
Reflect & 15 min & Record what changed in your understanding and what remains uncertain. \\
\bottomrule
\end{tabularx}

\subsection{Definition of Done for Each Week}
A weekly module is complete only when:
\begin{itemize}
  \item the assigned reading or free equivalent is completed;
  \item the week's artifact is committed or otherwise versioned;
  \item terminology and diagram notation are defined;
  \item the artifact is reviewed from at least one stakeholder's point of view;
  \item the mastery check is at least 80\%; and
  \item open questions are recorded instead of silently guessed.
\end{itemize}

% ============================================================================
\section{Authoritative Source Basis}

\begin{enumerate}
  \item Carnegie Mellon University Software Engineering Institute. \emph{Documenting Software Architectures}. \href{https://www.sei.cmu.edu/training/documenting-software-architectures/}{Official SEI course page}.
  \item Clements, P., Bachmann, F., Bass, L., Garlan, D., Ivers, J., Little, R., Merson, P., Nord, R., and Stafford, J. \emph{Documenting Software Architectures: Views and Beyond, 2nd Edition}. Addison-Wesley Professional, 2010. \href{https://www.sei.cmu.edu/library/documenting-software-architectures-views-and-beyond-second-edition/}{SEI book record}.
  \item Carnegie Mellon University Software Engineering Institute. \emph{Views and Beyond Collection}. \href{https://www.sei.cmu.edu/library/views-and-beyond-collection/}{SEI collection}.
  \item Nord, R., Clements, P., Emery, D., and Hilliard, R. \emph{A Structured Approach for Reviewing Architecture Documentation}, CMU/SEI-2009-TN-030. \href{https://www.sei.cmu.edu/library/a-structured-approach-for-reviewing-architecture-documentation/}{SEI report}.
  \item ISO/IEC/IEEE 42010:2022. \emph{Software, systems and enterprise --- Architecture description}, Edition 2. \href{https://www.iso.org/standard/74393.html}{Official ISO record}.
  \item Object Management Group. \emph{Unified Modeling Language (UML), Version 2.5.1}. \href{https://www.omg.org/spec/UML/}{Official OMG specification}.
\end{enumerate}

\vfill
\begin{tcolorbox}[
  enhanced,
  colback=Navy,
  colframe=Navy,
  boxrule=0pt,
  arc=2pt,
  left=12pt,right=12pt,top=10pt,bottom=10pt]
{\color{white}\bfseries Final objective}\par
{\color{white!90}Be able to take an unfamiliar software repository and systematically derive a stakeholder-driven, multi-view, ISO-42010-aligned architecture description that is understandable, reviewable, and useful to engineering, security, operations, and governance stakeholders.}
\end{tcolorbox}

\end{document}
